\section{Interfaces as Contracts}
\label{sec:contracts}

Every layer in \cref{sec:layers} makes a claim about its output under
assumptions about its input.  An assume--guarantee contract separates an
environment assumption $A$ from a component guarantee $G$
\cite{benveniste2007contracts}.  A contract theory may also define satisfaction,
compatibility, refinement, composition, and quotient for a chosen behavioral
model.  Robotics uses this idea at several levels: a mission contract can
describe temporal goals and contexts; a component contract can specify software
behavior; a dynamical contract can constrain input and output trajectories; and
a runtime-assurance specification can govern transfer of control authority.

Here, a \emph{contract} states obligations, \emph{verification} checks a stated
property, and \emph{assurance} is the broader argument and evidence supporting
confidence in system behavior.  An interface schema is not automatically a
contract.  Names, types, timestamps, and version fields become contractual only
when they participate in explicit assumptions or guarantees, are checked by a
defined validator, and have a specified consequence when the check fails.

\subsection{Mission Contracts}

CROME represents contextual robotic missions as hierarchical assume--guarantee
contracts and synthesizes sub-mission controllers
\cite{mallozzi2020crome,mallozzi2023contextual}.  Contract algebra can also
refine a mission against a component library or compute a repair when the
original specification is unrealizable \cite{mallozzi2022repair}.  These methods
make goals, contexts, refinement, and realizability explicit at the mission
boundary.  They focus on the mission specification; data age, provenance, and
implementation identity remain outside their scope.

\subsection{Component Contracts}

Luckcuck \emph{et al.} specify ROS nodes with first-order-logic contracts,
derive system properties with composition rules, and translate a ROS Contract
Language into executable monitors \cite{luckcuck2022compositional}.  This work
connects behavioral contracts to node implementations and demonstrates that
component assumptions can be made executable.  Its primary object is node
behavior; cross-robot state age, team membership, and heterogeneous deployment
compatibility are not the central object of the calculus.

\subsection{Dynamical and Multi-Agent Contracts}

Contracts for dynamical systems place assumptions on input trajectories and
guarantees on outputs.  Sharf \emph{et al.} provide satisfaction and refinement
tools for linear contracts and illustrate them with two vehicles
\cite{sharf2021dynamical}.  Chen \emph{et al.} use parameterized temporal-logic
contracts to compute invariant sets compositionally in networked systems
\cite{chen2019invariance}.  Liu \emph{et al.} synthesize decentralized
controllers for interconnected nonlinear systems from signal-temporal-logic
contracts \cite{liu2022stlcontracts}.

Assume--guarantee reinforcement learning instead represents neighboring-agent
behavior through regular-language abstractions and combines local satisfaction
probabilities into a global lower bound \cite{kazemi2024agrl}.  These approaches
directly address interaction and decomposition, but typically assume that the
declared neighbor abstraction, labels, communication observations, dynamics,
and timing semantics are available as modeled.
